\section{Technical validation}\label{valid}

Technical validation separates source reconciliation and measurement
checks from structural constraints, executable queries and the
operational ELK check. The evidence bundles identify the input
snapshot and checks executed.

\subsection{Assay validation}

The PowerSoil Pro extraction-control profile contained 1,165 ASVs,
with 84\% of reads assigned to ASVs shared with biological profiles.
Its nearest canonical profile was 46Dr1, with Bray--Curtis
dissimilarity 0.39. We excluded this control from the contaminant
screen. A library-order test compared nearby library numbers with
different-site pairs matched by geographic-distance decile and
compartment, using 999 library-order permutations. The Trip~1 and
Trip~3 Bray--Curtis contrasts were $-0.003$ ($p=0.15$) and
$+0.007$ ($p=1.0$), respectively. These contrasts summarize the
tested preparation-order association.

For XRF, the source audit reproduces 13,390 Trips~1--4 and 4,293
Trip~5 reported cells under their stated selection rules.
Repeated field sessions were reduced to site medians and paired
with the Trip~5 subsurface laboratory observation at 58 sites.
Among 12 analytes with positive values in both workflows at more
than 10 paired sites, the median Spearman correlation was 0.21
(range $-0.40$ to 0.48;
\path{evidence/xrf_audit/xrf_field_lab_agreement.tsv}).
This diagnostic describes site-level rank agreement among the
retained positive pairs.

\subsection{Competency questions and SPARQL queries}

Competency questions define the intended retrieval tasks
(Table~\ref{tab:cq}). The archived query workload covers site,
geochemical, molecular-provenance and taxonomic data. The
Site~10 XRF query has an archived 46-row result. Source-derived
expected result tuples also specify the complete 48-observation
genus profile for run \texttt{ERR16061083},
360-row environment--taxon join, 70 site coordinates and six
PCR no-template control roles. The release validator compares full
result multisets, preserving duplicate rows and requiring exact
identifiers. Numeric comparisons use absolute and relative
tolerances of $10^{-12}$; coordinate comparisons use an absolute
tolerance of $10^{-10}$ degrees. The 26 checks cover the printed
queries across direct and application-mediated transports, default
and explicit datasets, and a 20-row taxonomic API result.

Six additional checks verify the advertised asserted-only dataset,
the rejection of unavailable query modes, the absence of unreleased
graphs, the default graph cardinality and all 712 pH
process--specimen links against the reconciled source module.
Fourteen protocol checks cover dataset selection, read-only access
and preservation of RDF terms across the supported graph and
SELECT serializations. Browser validation executes all nine portal
examples and checks the three versioned download links. Full HTTP
downloads verify artifact sizes and checksums; the complete asserted
N-Quads export is parsed and checked for its graph context and triple
count. All checks passed on the candidate and public service. The
\href{https://rubalkhali.science/downloads/kg/3.0.0/post-deployment-validation.json}{public validation report}
binds these results to the release manifest.
Separate structural tests cover additional module types.

\begin{table}[h]
\centering
\caption{Competency questions defining the retrieval workload.}
\label{tab:cq}
\small
\begin{tabular}{p{0.08\textwidth} p{0.84\textwidth}}
\toprule
\textbf{ID} & \textbf{Competency question} \\
\midrule
CQ1 & What are the XRF analyte concentrations measured at a given sampling site? \\
CQ2 & Which sampling sites are recorded, and what are their GPS coordinates? \\
CQ3 & Which biome and environmental features are assigned to each sampling site? \\
CQ4 & What are the instrument-reported Light Elements values across the XRF measurements? \\
CQ5 & Which DNA extraction records (extract, concentration, protocol) exist for a given sample? \\
CQ6 & What is the full provenance chain from a soil sample to its sequencing run? \\
CQ7 & What are the complete genus-level read counts and relative read abundance for run \texttt{ERR16061083}? \\
\bottomrule
\end{tabular}
\end{table}

We show several example SPARQL queries. Listing~\ref{lst:sparql_xrf}
retrieves XRF measurements for a specific site, requiring traversal of
the measurement process, quality, and value entities.

\begin{lstlisting}[caption={SPARQL query to retrieve field XRF
    measurements for Site 10. The query starts from the field-XRF
    measuring processes whose target is the site and walks to the
    individual analyte concentrations.},
  label={lst:sparql_xrf}, basicstyle=\ttfamily\scriptsize,
  frame=single]
PREFIX rak: <https://rubalkhali.science/kb/RAK_>
PREFIX sio: <http://semanticscience.org/resource/SIO_>
PREFIX rdfs: <http://www.w3.org/2000/01/rdf-schema#>

SELECT ?processLabel ?siteLabel ?analyte ?concentration
WHERE {
  ?process a rak:0000025 ;
           rdfs:label ?processLabel ;
           sio:000291 ?site ;
           sio:000229 ?value .
  FILTER(CONTAINS(LCASE(STR(?processLabel)), "field xrf"))
  ?site rdfs:label ?siteLabel .
  FILTER(?siteLabel = "Site 10")
  ?value sio:000215 ?quality ;
         rak:2000012 ?concentration .
  ?quality a ?qualityClass .
  ?qualityClass rdfs:label ?analyte .
  FILTER(?qualityClass != rak:0000032)
}
ORDER BY ?processLabel ?analyte
\end{lstlisting}

The query selects field-XRF processes targeting Site~10, then
retrieves each output and its analyte quality. It returns one row
per process/analyte combination, excluding the aggregate Light
Elements channel. The concentration units are undocumented.

The release workflow executes this exact query
(\path{sparql/field_xrf_site10.rq}) against the checksummed base, site
and XRF modules and requires the fixed expected cardinality of
46 analyte-value rows; it returned exactly 46 rows for Site~10,
24 from field process Test~5847 and 22 from Test~5848, and the result
table and validation record are archived under
\path{evidence/competency-query/}.
This execution checks the archived query's cardinality; release
validation also requires comparison of stable identifiers and values
with the expected table.

Listing~\ref{lst:sparql_taxon} retrieves the complete genus profile for
run \texttt{ERR16061083},
joining across bioinformatic workflows, FASTQ datasets, and taxon
observations.

\begin{lstlisting}[caption={Complete genus-level profile for run ERR16061083, ordered by descending read count and lineage.},label={lst:sparql_taxon},basicstyle=\ttfamily\scriptsize,frame=single]
PREFIX rak: <https://rubalkhali.science/kb/RAK_>
PREFIX sio: <http://semanticscience.org/resource/SIO_>
PREFIX rdfs: <http://www.w3.org/2000/01/rdf-schema#>
PREFIX xsd: <http://www.w3.org/2001/XMLSchema#>

SELECT ?proc ?fastq ?runLabel ?lineage ?taxon ?taxonLabel
       ?count ?relativeAbundance
WHERE {
  VALUES ?runLabel { "FASTQ dataset for ERR16061083" }
  ?proc a rak:0000071 ;
        sio:000230 ?fastq ;
        sio:000229 ?dsAbs, ?dsRel .
  ?fastq rdfs:label ?runLabel .
  ?dsAbs a rak:0000074 ; rdfs:label ?absLabel ;
         sio:000059 ?valAbs .
  ?dsRel a rak:0000075 ; rdfs:label ?relLabel ;
         sio:000059 ?valRel .
  FILTER(STRENDS(STR(?absLabel), "(Genus)"))
  FILTER(STRENDS(STR(?relLabel), "(Genus)"))
  ?valAbs rak:2000026 ?count ;
          rak:2000025 ?lineage ;
          sio:000215 ?qualAbs .
  ?valRel rak:2000020 ?relativeAbundance ;
          rak:2000025 ?lineage ;
          sio:000215 ?qualRel .
  ?qualAbs sio:000011 ?taxon, ?fastq .
  ?qualRel sio:000011 ?taxon, ?fastq .
  ?taxon rdfs:label ?taxonLabel .
  FILTER(?taxon != ?fastq)
}
ORDER BY ?runLabel DESC(xsd:double(?count)) ?lineage ?proc
\end{lstlisting}

The query binds both count and relative-abundance values to the same
workflow, source FASTQ dataset, taxon and classifier lineage. Dataset
labels currently encode rank through the suffix \texttt{(Genus)}.
The query uses that emitted convention explicitly and returns all
48 observations for the named run. The source-derived expected table
contains the complete profile, including each classifier lineage.
The separate \path{scripts/release/kg/taxonomy_all_runs.rq} retrieves
all runs from a locally loaded snapshot. That unrestricted export
is a long-running operation. The environmental example below additionally binds
field temperature and specimen collection to the same site visit.

\begin{lstlisting}[caption={SPARQL query to retrieve Pseudomonas relative abundance at sites with high temperature.}, label={lst:sparql_env_tax}, basicstyle=\ttfamily\scriptsize, frame=single]
PREFIX rak: <https://rubalkhali.science/kb/RAK_>
PREFIX sio: <http://semanticscience.org/resource/SIO_>
PREFIX pato: <http://purl.obolibrary.org/obo/PATO_>
PREFIX ncbitaxon: <http://purl.obolibrary.org/obo/NCBITaxon_>
PREFIX rdfs: <http://www.w3.org/2000/01/rdf-schema#>
SELECT ?visit ?visitLabel ?site ?siteLabel ?sample ?fastq
       ?temperatureProcess ?tempVal ?temp ?lineage ?relAbundance
WHERE {
  ?visit a rak:0000003 ; rdfs:label ?visitLabel ;
         sio:000291 ?site .
  ?site rdfs:label ?siteLabel .
  ?temperatureProcess sio:000068 ?visit ;
                      sio:000229 ?tempVal .
  ?tempVal rak:2000003 ?temp ; sio:000215 ?tempQual .
  ?tempQual a pato:0000146 ; sio:000011 ?site .
  FILTER(?temp > 35.0)
  ?sampling a rak:0000024 ; sio:000068 ?visit ;
            sio:000291 ?site ; sio:000229 ?collection .
  ?collection sio:000059 ?sample .
  ?ext sio:000230 ?sample ; sio:000229 ?dna .
  ?libPrep a rak:0000065 ; sio:000230 ?dna ;
           sio:000229 ?lib .
  ?seq a rak:0000066 ; sio:000230 ?lib ;
       sio:000229 ?fastq .
  ?bioinf a rak:0000071 ; sio:000230 ?fastq ;
          sio:000229 ?dsRel .
  ?dsRel a rak:0000075 ; rdfs:label ?datasetLabel ;
         sio:000059 ?valRel .
  FILTER(STRENDS(STR(?datasetLabel), "(Genus)"))
  ?valRel rak:2000020 ?relAbundance ;
          rak:2000025 ?lineage ; sio:000215 ?qualRel .
  ?qualRel sio:000011 ncbitaxon:286, ?fastq .
}
ORDER BY ?visit ?sample ?fastq ?tempVal ?lineage
\end{lstlisting}

Each result retains the collection visit, specimen, sequence dataset,
temperature observation and genus lineage. A site's visits are joined
through their process provenance. This provides an environmental
temperature associated with the collection visit of each returned
microbial profile.

Regression fixtures execute both queries on deliberately confusable
inputs. The taxonomy fixture assigns the same genus different
abundances in two profiles and adds a species-rank observation; the
query returns exactly the two matching genus observations. The visit
fixture gives one site two visits at 40 and 20\degree C; the
temperature-filtered query returns only the specimen collected during
the 40\degree C visit. These tests check the implemented joins;
release-level validation additionally requires expected result tables
for the complete modules.

\subsection{Source reconciliation and measurement validation}

The archived evidence workflow produced a row-level sample ledger and
reconciliation summaries for specimens, extracts, libraries, runs and
community profiles. The 2,516 eligible source rows reconcile exactly
to 2,516 generated sample individuals; the 1,647 extraction-eligible
rows reconcile to 1,647 DNA extracts; and 1,242 resolved sequence rows
reconcile to 1,242 amplicon libraries and 1,242 FASTQ datasets. A
site-alias gate independently matches each Trip~1 numeric identifier
61--64 to a named catalogue individual by exact coordinates and
rejects coordinate or IRI disagreement, accounting for all 36 affected
source rows.

The environmental-metadata gate reproduced 274 nonblank source records
and verified all 24 exact-cell correction-ledger dispositions against
the corresponding raw cells. All retained field temperatures,
pressures and relative humidities fell within their declared physical
ranges; the unresolved Trip~5 site~40 value of 194\% was preserved in
the ledger but is absent from both the curated table and the
environmental RDF module. The XRF provenance audit reproduced all
13,390 reported Trips~1--4 cells under the maximum-positive rule and
all 4,293 reported Trip~5 cells under the last-reported rule, with no
mismatched reported cells, and verified that the canonical laboratory
input comprises all 725 processed records. An independent
chemical-identifier gate checked all 93 instrument channels against
the pinned ChEBI release 247 and the checksummed PubChem snapshot of
23 July 2026 (Section~\ref{methods:xrf}); all
rows passed the declared rules, with explicit nulls for the aggregate
LE channel and unsupported oxides. The pH reconciliation accounts for 1,168 workbook rows and
reproduces the five admission dispositions in
Section~\ref{methods:ph}. The string-based join associates each
admitted row with a specimen identifier; the Trip~4 physical-replicate
correction requires the separate provenance reconciliation described
there.

\subsection{ShEx schema validation}

The archived ShEx run used the shapes described in
Section~\ref{subsec:schema_def}. Its tractable modules passed their
assigned shapes. The pH module was
validated separately over all 712 measurement processes, values and
acidity qualities and all 29 sessions, and the laboratory-control
module passed its shape and scientific-invariant checks. Four negative
fixtures each alter one constraint: missing site coordinates, an invalid
value datatype, reversed value-to-quality direction, or an XRF process
with neither input nor target. The released shapes reject all four
fixtures. The
multi-gigabyte taxonomy ABox remains outside this memory-bounded run;
it is instead parsed in full by an independent Turtle parser and
scanned by a streaming structural validator, which accounted for
44,528,482 triples in 5,953,690 subject blocks with zero structural
violations. This checks full-file syntax and the generator's declared structural
relationships. The separate listing validator checks printed Turtle
against the modules available for that run. Taxonomy and external
reference checks require the corresponding pinned payloads.

\subsection{Logical consistency}

The 9 September validation loaded 781,044 asserted axioms from eleven
ontology IRIs and ELK returned \texttt{true} for its consistency check.
The checked union included the ecosystem, pH and laboratory-control
modules. It excluded the taxonomy ABox and large reference modules;
imported axioms were not added to the asserted union, and class-hierarchy
precomputation was omitted. An import-order regression confirmed that
modules loaded through imports and then encountered explicitly retain
their asserted axioms exactly once.

The generated assertions include \texttt{xsd:double} literals, placing
the asserted graph outside OWL~2~EL's datatype restrictions
\cite{owl2profiles}. An OWLAPI profile diagnostic on the tractable
asserted core found 40,482 illegal-datarange occurrences, including
40,481 uses of \texttt{xsd:double}. Diagnostics caused by missing
import declarations were recorded separately. ELK's operational
result describes the asserted union under the constructs processed
by this ELK configuration.

Validation of the separate finite-rule implementation uses synthetic examples for long
hierarchies, transitive cycles, both project property-chain lengths,
restriction and intersection classification, and PCR-control ancestry.
The tests compare bounded, class-partitioned execution with an
unbounded reference closure and check recovery from an interrupted
run. Release-specific checks compare the complete control-role
inventory with its source, retrieve consequences of the measurement
chains and require the inferred graph to be disjoint from the
asserted graph. Completion requires the zero-addition pass specified
in Section~\ref{subsec:owl_rdf}, followed by the explicit clash
diagnostics.

The separate label gate merged 800,026 triples and confirmed exactly
one label for each of 106,833 labelled \texttt{RAK\_} subjects.
Current logs, input hashes and validator sources are preserved under
\path{evidence/final-validation-20260909/}; the earlier run remains
under \path{evidence/semantic-validation/}.

\subsection{Compute-time metrics}\label{sec:compute_metrics}

Table~\ref{tab:compute} preserves development timings for a
nine-ABox snapshot. The original benchmark lacks the complete
hardware, repetition and cache information needed for a portable
performance comparison. It differs from the current module set.
The query timings range from 5~ms to a timeout above 300~s.

\begin{table}[h]
\centering
\caption{Historical development timings; these describe an earlier
  nine-ABox snapshot.}
\label{tab:compute}
\small
\begin{tabular}{lr}
\toprule
\textbf{Step} & \textbf{Time} \\
\midrule
RDF module loading (Virtuoso, 9 ABox files)  & 140\,s \\
Selected rule expansion (3 passes)           & 114\,min \\
ELK check (historical benchmark)           & 1.6\,s \\
CQ1 -- XRF measurements, Site 10            & 23\,ms \\
CQ2 -- Sampling sites + GPS                 & 5\,ms \\
CQ3 -- Biomes per site                      & 5\,ms \\
CQ4 -- Light Elements reported values       & 8\,ms \\
CQ5 -- DNA extraction records               & 1{,}342\,ms \\
CQ6 -- Provenance chain                     & 255\,ms \\
CQ7 -- Taxonomic composition                & $>$300\,s \\
\bottomrule
\end{tabular}
\end{table}

CQ7 exceeded the five-minute timeout in that benchmark. A subsequent
global 100-row preview and an equivalent two-branch formulation also
exceeded 300~s on the idle release candidate. These runs are archived
with their exact queries. The current CQ7 demonstrates a complete
profile for a named run; its acceptance test compares all observations
with an independently derived source table. Performance measurements
must identify the query and graph snapshot used.

\subsection{Portal access}

The portal (\url{https://rubalkhali.science/}) retrieves data through
backend SPARQL queries and renders site and specimen views. This is
an access interface. Independent validation compares query results
with source-derived expected values; a user study would additionally
measure task completion, errors and usability.

\subsection{Identifiers and reuse metadata}

Each named entity has a project-domain IRI. The manifest identifies
payload files, and the data dictionary specifies curated fields,
units and missing-value conventions. RDF provenance links travel
with the modules and support reuse through SIO, ENVO, ChEBI, PATO,
UO and validated taxonomic identifiers \cite{fair}.

The public endpoint provides read access to the deployed graph.
A fixed archival package additionally requires a versioned deposit,
verified payload access, a persistence policy and explicit licences
for project data, code and externally sourced material. Data
Availability identifies the remaining sequence-release gate.
